\documentclass[11pt]{article}
\usepackage[a4paper,margin=28mm]{geometry}
\usepackage{amsmath,amssymb,amsthm,hyperref}
\newtheorem{theorem}{Theorem}
\title{Excluding the first four odd-degree parity-bound candidates}
\author{Mathematical proof independently reviewed; not formally verified}
\date{30 September 2026}
\begin{document}
\maketitle
Write $v=v_2$. We use the parity-count and resonance conclusions proved
in \path{../independent_directions/optimized_two_adic_filter.tex}.
For odd degree $r$, if a rational equal-weight rule has
$2^{r-1}\le K<2^r$, and $N=2^a$ is the least power of two making all
nodes $y_i=Nx_i$ integral in $\mathbb Z_2$, then exactly $2^{r-1}$
normalized nodes are odd, $N\mid r+1$, and
\[
 v(K)=(r+1)/N-1.\tag{1}
\]
The case where all original nodes are integral forces $2^r\mid K$.

\begin{theorem}
For $r=2^a-1$ with $a\ge3$, no rational equal-weight uniform-interval
rule of degree $r$ has $K=2^{r-1}+1$ nodes. Consequently
\[
 R_r\ge2^{r-1}+2.
\]
In particular $R_7\ge66$.
The nodes need not be in $[0,1]$ for this lower bound.
\end{theorem}
\begin{proof}
Suppose such a rule exists. Its $K$ is odd, so (1) forces
$N=r+1=2^a$. Exactly $2^{r-1}$ normalized nodes are odd, and there is
one even node $z$. Put
\[
 k=(r-1)/2=2^{a-1}-1,\qquad
 f(y)=y(y^2-1)^k,\qquad g(y)=(y^2-1)^k.
\]
Both polynomials have degree at most $r$. For odd $y\in\mathbb Z_2$,
$y^2-1\in8\mathbb Z_2$, so $f(y),g(y)\in2^{3k}\mathbb Z_2$.

The integral of $f(Nx)$ is
\[
 A=\sum_{j=0}^k {k\choose j}(-1)^{k-j}
           \frac{N^{2j+1}}{2j+2}.
\]
The $j=0$ term has valuation $a-1$. For $j\ge1$, the valuation is
at least $(2j+1)a-1-v(j+1)>a-1$. Thus $v(KA)=a-1$.
Exactness gives $f(z)\equiv KA\pmod{2^{3k}}$, and $3k>a-1$.
Since $z$ is even, $v(f(z))=v(z)$, proving
\[
 v(z)=a-1.\tag{2}
\]

As $k$ is odd, expansion of $g(z)-(-1)^k$ has leading nonzero term
$k(-1)^{k-1}z^2$. All later terms have greater valuation, and (2) yields
\[
 v\bigl(g(z)-(-1)^k\bigr)=2a-2.\tag{3}
\]
On the other hand, writing $B=K\int_0^1g(Nx)\,dx$, we have
\[
 B-(-1)^k=(-1)^k(K-1)
  +K\sum_{j=1}^k {k\choose j}(-1)^{k-j}
            \frac{N^{2j}}{2j+1}.
\]
Here $v(K-1)=r-1=2^a-2\ge2a$, and every summand in the sum has
valuation at least $2a$. Therefore $v(B-(-1)^k)\ge2a$.
Exactness for $g$ gives $g(z)\equiv B\pmod{2^{3k}}$; since
$3k\ge2a$, this contradicts (3).
\end{proof}

\begin{theorem}[Four consecutive exclusions]
For $r=2^a-1$, $a\ge3$, one has
\[
 R_r\ge2^{r-1}+5.
\]
In particular $R_7\ge69$.
\end{theorem}
\begin{proof}
The preceding theorem excludes the first candidate. Suppose now that
$K=U+s$, where $U=2^{r-1}$ and $s\in\{2,3,4\}$.
The resonance condition (1) immediately excludes $s=4$: it would give
$v(K)=2$ and $(r+1)/N=3$, impossible because $r+1$ and $N$ are powers
of two.

Retain $k=(r-1)/2$, $f(y)=y(y^2-1)^k$, and $g(y)=(y^2-1)^k$.
For $s=2$, (1) gives $N=2^{a-1}$ and two even normalized nodes; for
$s=3$, it gives $N=2^a$ and three even normalized nodes.
In both cases the same uniquely minimal linear term in the integral
of $f(Nx)$ gives
\[
 v\left(K\int_0^1f(Nx)\,dx\right)=a-1.\tag{4}
\]
All odd-node contributions are in $2^{3k}\mathbb Z_2$.

We use an elementary valuation observation. For even $z$ put
\[
 h(z)=g(z)+1=z^2Q(z^2),\qquad Q(0)=k\ \text{odd},
\]
where $Q$ has integer coefficients. For all $z$ of a fixed finite
valuation $b\ge1$, the odd quantity $h(z)/2^{2b}$ has the same residue
$Q(2^{2b})$ modulo eight: write $z=2^bu$ and use $u^2\equiv1\pmod8$.
For a sum of two or three such $h(z)$, let $b$ be the smallest valuation
of the nonzero $z$'s. If one or three terms attain $b$, the sum has
valuation $2b$. If exactly two terms attain $b$, their normalized sum
is two or six modulo eight. A possible third term contributes zero or
four modulo eight, so the sum has valuation $2b+1$.

For $s=2$, expansion of the integral gives
\[
 K\int_0^1g(Nx)\,dx+2
 =-U+K\sum_{j=1}^k {k\choose j}(-1)^{k-j}
                 \frac{N^{2j}}{2j+1}.
\]
Its $j=1$ term has valuation $1+2(a-1)=2a-1$; all other terms,
including $U$, have greater valuation. Exactness and $3k>2a-1$ show
that the sum of the two $h(z)$ has valuation $2a-1$.
The observation forces both even nodes to have valuation $a-1$.
Their $f$ values, divided by $2^{a-1}$, are both odd, so their sum
has valuation at least $a$, contradicting (4).

For $s=3$ and $a\ge4$, the analogous integral plus three has exact
valuation $2a$: its $j=1$ term is uniquely minimal, since
$v(U)=2^a-2>2a$. The observation forces every even node to have
valuation at least $a$. Their $f$ values then sum to an element of
$2^a\mathbb Z_2$, again contradicting (4).
For the remaining case $a=3,s=3$, the first two terms combine as
$-64+67\cdot64=64\cdot66$, with exact valuation seven; all later
terms have valuation at least twelve. Hence the sum of the three
$h(z)$ has valuation seven. The observation forces exactly two nodes
of valuation three and the remaining node of valuation at least four.
Their $f$ values sum with valuation at least four, contradicting the
required valuation two in (4).
\end{proof}

\paragraph{Review and origin of the strengthening.}
The initial one-node exclusion and the degree-seven congruences below
were checked independently by two parallel reviewers. One reviewer
then supplied the two-node exclusion and the general $s=2,3,4$
extension; the three-node degree-seven case was independently derived.
The resulting four-exclusion proof was checked separately case by case.
This is mathematical peer review, not formal proof verification.

\paragraph{The degree-seven calculation directly.}
If $K=65$, the normalization is $N=8$, with $64$ odd nodes and one
even node $z$. The polynomial $y(y^2-1)^3$ takes values in
$512\mathbb Z_2$ on odd nodes, whereas
\[
 65\int_0^1 8x((8x)^2-1)^3\,dx
 =65(2^{18}-2^{14}+384-4)
\]
has valuation two. Thus $v(z)=2$. But $(y^2-1)^2$ takes values in
$64\mathbb Z_2$ on odd nodes and
\[
 65\int_0^1((8x)^2-1)^2\,dx
 =65\left(\frac{4096}{5}-\frac{128}{3}+1\right)
 \equiv1\pmod{64}.
\]
For $v(z)=2$, $(z^2-1)^2\equiv33\pmod{64}$, a contradiction.

\section*{A further obstruction for symmetric degree-seven rules}
The general lower bound is $R_7\ge69$. The next candidate cannot be
attained by a rule symmetric about $1/2$.
\begin{theorem}
No symmetric rational equal-weight degree-seven rule has $69$ nodes.
\end{theorem}
\begin{proof}
The resonance condition gives least normalization $N=8$, with $64$
odd nodes. Center the normalized coordinates as $u=8x-4$. Symmetry
makes the five even coordinates a multiset $0,v,-v,w,-w$, where
$v,w\in2\mathbb Z_2$; initially either may be zero. Put
$g(u)=(u^2-1)^3$ and $h(u)=g(u)+1$. The odd coordinates contribute
$64$ modulo $512$ to the sum of $h$. On the other hand,
\[
69\int_{-1}^1h(4t)\,\frac{dt}{2}-64
=69\left(\frac{4096}{7}-\frac{768}{5}+16\right)-64
=\frac{16\cdot67411}{35},
\]
which has valuation $4$. Thus $v_2(h(v)+h(w))=3$. The fixed-valuation
residue argument used above shows that two nonzero even inputs give
valuation $2b$ if their valuations differ, and $2b+1$ if both have
valuation $b$. Consequently $v_2(v)=v_2(w)=1$.

It follows that $v^2,w^2\equiv4\pmod{32}$, hence
$g(v),g(w)\equiv27\pmod{32}$. The even-node contribution to $g$ is
therefore
\[
 g(0)+2g(v)+2g(w)\equiv-1+4\cdot27=43\pmod{64}.
\]
The odd-node contribution vanishes modulo $512$, whereas exactness
requires
\[
69\int_{-1}^1g(4t)\,\frac{dt}{2}
=69\left(\frac{4096}{7}-\frac{768}{5}+15\right)
\equiv11\pmod{64},
\]
a contradiction. This argument does not exclude nonsymmetric
$69$-node rules.
\end{proof}


\section*{An explicit seventy-node rule}
There is a symmetric rational degree-seven rule with $70$ nodes.
Take $x=(18\pm s)/36$, with the following multiplicities for each sign:
\[
\begin{array}{c|rrrrrrrrrrr}
s&1&3&4&5&7&10&11&12&13&15&17\\\hline
c_s&3&3&1&3&9&1&2&1&4&4&4.
\end{array}
\]
Indeed $\sum c_s=35$, and
\[
 \sum c_ss^2=3780=\frac{35\cdot18^2}{3},\qquad
 \sum c_ss^4=734832=\frac{35\cdot18^4}{5},\qquad
 \sum c_ss^6=170061120=\frac{35\cdot18^6}{7}.
\]
Symmetry supplies the odd centered moments. The rule therefore has all
required moments through degree seven. Repeated nodes are allowed in
an equal-weight rule. The standalone verifier
\path{verify_degree7_seventy.py} checks all eight ordinary moments with
exact rational arithmetic.

Consequently
\[
 \boxed{69\le R_7\le70},
\]
and the minimum number of nodes among symmetric rational degree-seven
rules is exactly $70$. Whether a nonsymmetric $69$-node rule exists
remains open in these notes.

\end{document}
